\subsection{Statistical Experiments and Likelihoods}

A statistical problem starts with observable data and a family of probability laws that could have generated it.  The model is not the data itself; it is the mathematical language in which uncertainty about the data-generating mechanism is expressed.

\begin{definition} \label{def:statistical-experiment-and-model} \hlt{Statistical Experiment and Model}\\
A \hlt{statistical experiment} is a measurable space $(\mathcal{X},\mathcal{A})$ together with a family of probability measures
\[
\mathcal{P}=\{P_\theta:\theta\in\Theta\}
\]
on $(\mathcal{X},\mathcal{A})$.  The set $\Theta$ is the \hlt{parameter space}; a realization $X\in\mathcal{X}$ is the \hlt{sample}; and $\mathcal{P}$ is the \hlt{statistical model}.  If the true law belongs to $\mathcal{P}$, the model is \hlt{well specified}; otherwise it is \hlt{misspecified}.
\end{definition}

\begin{definition} \label{def:dominated-model-likelihood-score-information} \hlt{Dominated Models, Likelihood, Score, and Information}\\
The model is \hlt{dominated} by a measure $\mu$ if $P_\theta\ll\mu$ for every $\theta\in\Theta$.  Writing $p_\theta=dP_\theta/d\mu$, the \hlt{likelihood} at observation $x$ is
\[
L_x(\theta)=p_\theta(x),
\qquad
\ell_x(\theta)=\log p_\theta(x)
\]
when $p_\theta(x)>0$.  If $\Theta\subseteq\R^d$ and differentiation is justified, the \hlt{score} and \hlt{Fisher information} are
\[
S_\theta(X)=\nabla_\theta\ell_X(\theta),
\qquad
I(\theta)=\mathbb{E}_\theta\!\left[S_\theta(X)S_\theta(X)^\top\right].
\]
\end{definition}

\begin{example} \label{ex:bernoulli-and-normal-likelihoods} \hlt{Bernoulli and Normal Likelihoods}
\begin{enumerate}[label=(\roman*)]
\setlength{\itemsep}{0pt}
\item If $X_1,\ldots,X_n$ are independent Bernoulli$(p)$ observations, then
\[
L_x(p)=\prod_{i=1}^n p^{x_i}(1-p)^{1-x_i}
=p^{\sum_i x_i}(1-p)^{n-\sum_i x_i}.
\]
The data enter the likelihood only through the count $\sum_iX_i$.
\item If $X_1,\ldots,X_n$ are independent $\mathcal{N}(\mu,\sigma^2)$ observations with known $\sigma^2$, then
\[
\ell_x(\mu)=-\frac{n}{2}\log(2\pi\sigma^2)-\frac{1}{2\sigma^2}\sum_{i=1}^n(x_i-\mu)^2.
\]
The maximizer is the sample mean $\overline{X}$.
\end{enumerate}
\end{example}

\begin{definition} \label{def:identifiability-and-parametric-functionals} \hlt{Identifiability and Functionals}
\begin{enumerate}[label=(\roman*)]
\setlength{\itemsep}{0pt}
\item The parameterization is \hlt{identifiable} if $P_\theta=P_{\theta'}$ implies $\theta=\theta'$.
\item A \hlt{statistical functional} is a map $T:\mathcal{P}\to\mathcal{T}$, such as $T(P)=\int x\,dP(x)$ or $T(P)=\operatorname{median}(P)$.
\item In a parametric model, a \hlt{target parameter} may be a function $\psi(\theta)$ rather than the full parameter $\theta$.
\end{enumerate}
\hlt{Key Insight}: Estimation should target a law or functional that the experiment can actually distinguish.
\end{definition}

\begin{remark} \label{rem:likelihood-is-relative-evidence} \hlt{Likelihood as Relative Evidence}\\
For fixed data $x$, the likelihood is a function of $\theta$, not a probability distribution on $\Theta$.  Ratios $L_x(\theta_1)/L_x(\theta_0)$ compare how well two model laws explain the same observation.
\end{remark}
